\documentclass[12pt,a4wide]{report}

\usepackage{amsthm,amssymb,mathrsfs,setspace,pstcol}
\usepackage{play}
\usepackage{epsfig}
\usepackage[nottoc]{tocbibind}
\usepackage{graphicx}
\usepackage{listings}
\usepackage{xcolor}
\usepackage{hyperref}
\usepackage{float}
\usepackage{adjustbox}

\renewcommand{\chaptermark}[1]{\markboth{#1}{}}
\renewcommand{\sectionmark}[1]{\markright{\thesection\ #1}}

\input xy
\xyoption{all}

\theoremstyle{plain}
\newtheorem{theorem}{Theorem}[section]
\newtheorem{lemma}[theorem]{Lemma}
\newtheorem{corollary}[theorem]{Corollary}
\newtheorem{proposition}[theorem]{Proposition}

\theoremstyle{definition}
\newtheorem{definition}[theorem]{Definition}
\newtheorem{example}[theorem]{Example}
\newtheorem{notation}[theorem]{Notation}

\theoremstyle{remark}
\newtheorem{remark}[theorem]{Remark}

% Code listing configuration
\lstset{
    language=Solidity,
    basicstyle=\ttfamily\small,
    keywordstyle=\color{blue}\bfseries,
    commentstyle=\color{green!60!black},
    stringstyle=\color{red},
    numbers=left,
    numberstyle=\tiny\color{gray},
    stepnumber=1,
    numbersep=5pt,
    backgroundcolor=\color{gray!10},
    frame=single,
    breaklines=true,
    breakatwhitespace=true,
    tabsize=2,
    showstringspaces=false
}

\renewcommand{\baselinestretch}{1.5}
\begin{document}

% --------------- Title page -----------------------
\begin{titlepage}
\enlargethispage{2cm}

\begin{center}

\vspace*{-1cm}

\textbf{\Large BLOCKCHAIN-BASED PUBLIC FUND TRACKING \& UTILIZATION SYSTEM}\\[10pt]
\vspace*{1.85 cm}
                          A Project Report Submitted \\
                     in Partial Fulfilment of the Requirements  \\
                     for the Degree of  \\
                          \vspace{4.10 mm}
                   {\Large \bf Bachelor of Technology } \\
                   in \\
                   {\large \bf Computer Science and Engineering }\\
                    with Specialization in\\
                   {\large \bf Cyber Security } \\
                      \vspace{3.85mm}
                   {\em  by} \\ \vspace{2.15mm}

{\large \bf [Student Name 1] - [Roll No. 1] }\\
{\large \bf [Student Name 2] - [Roll No. 2] }\\
{\large \bf [Student Name 3] - [Roll No. 3] }\\
{\large \bf [Student Name 4] - [Roll No. 4] }\\

\begin{figure}[h]
  \begin{center}
  \includegraphics[width=5cm, height=2.85cm]{logo2.jpg}
  \end{center}
\end{figure}
\vspace*{0.1cm}
{\em\large to }\\
{\bf \mbox{DEPARTMENT OF CSE-CYBER SECURITY}} \\
{\bf \mbox{INDIAN INSTITUTE OF INFORMATION TECHNOLOGY}}\\
{\bf KOTTAYAM-686635, INDIA}\\
{\it November 2024}

\end{center}
\end{titlepage}

\clearpage

% --------------- Abstract page -----------------------
\pagenumbering{roman} \setcounter{page}{2}
\begin{center}
{\Large{\bf{ABSTRACT}}}
\addcontentsline{toc}{chapter}{Abstract}
\end{center}

This project presents \textbf{PFTUS (Public Fund Tracking \& Utilization System)}, a blockchain-based solution addressing transparency, accountability, and corruption in public fund management. The system implements milestone-based escrow smart contracts that release funds only upon verified milestone completion.

Key innovations: milestone-based escrow, NFT approval stamps, zero-knowledge proofs, public trust index, and complete audit trail. The architecture consists of four smart contracts (PublicFundProject, ApprovalNFT, ImpactScoreCalculator, ZKVerifier) and a Next.js frontend with dashboards for Government, Contractor, Auditor, and Public users.

The implementation demonstrates blockchain's potential to revolutionize public fund management through real-time transparency, automated accountability, and immutable audit trails. Testing on Polygon Mumbai testnet and local Hardhat network shows promising results.

\clearpage

\tableofcontents
\clearpage
\listoffigures
\listoftables
\clearpage

\newpage

\pagenumbering{arabic}
\setcounter{page}{1}

% =========== Main chapters starts here ===========

\chapter{Introduction}

This project presents \textbf{PFTUS (Public Fund Tracking \& Utilization System)}, a blockchain-based platform that addresses critical challenges in public fund management through transparency, accountability, and immutability. Traditional centralized systems suffer from lack of transparency, delayed reporting, corruption risks, and inadequate audit trails \cite{blockchain-gov}.

Public fund management is a cornerstone of effective governance, involving billions of dollars annually in infrastructure projects, welfare schemes, and public services. However, traditional systems face significant challenges including opacity in fund allocation, delayed utilization reporting, corruption risks, and limited public access to information. These issues erode public trust and hinder effective governance.

\section{Background and Context}

Blockchain technology, introduced through Bitcoin \cite{nakamoto2008}, has evolved to enable decentralized applications through smart contracts \cite{buterin2014}. Smart contracts, first conceptualized by Nick Szabo \cite{szabo1994}, are self-executing contracts with terms directly written into code, enabling automated enforcement without intermediaries.

The application of blockchain in governance has shown promising results globally. Estonia's e-governance system \cite{estonia2017} demonstrates blockchain's potential for digital identity and secure record-keeping. Georgia's blockchain-based land registry \cite{georgia2016} provides transparent property records. These implementations highlight blockchain's ability to enhance transparency and reduce corruption in public administration.

\section{System Overview}

PFTUS leverages smart contracts to create an immutable, transparent system for tracking public fund allocation and utilization. The system enables: transparent fund tracking, automated milestone-based fund release, immutable audit trails, real-time monitoring, and public accountability.

The system architecture consists of four interconnected smart contracts deployed on blockchain networks, a web-based frontend application, and decentralized storage for documents. This combination ensures that all fund-related transactions are recorded immutably while maintaining user-friendly interfaces for all stakeholders.

\section{Key Stakeholders}

The system serves four primary stakeholders, each with distinct roles and responsibilities:

\begin{itemize}
    \item \textbf{Government:} Creates projects, defines milestones, allocates funds to smart contract escrow, and monitors project progress. Government users have administrative privileges to create and manage projects.
    
    \item \textbf{Contractors:} Execute project work, request milestone releases, submit completion proofs, and receive payments upon approval. Contractors can view their assigned projects and track milestone status.
    
    \item \textbf{Auditors:} Review submitted proofs, verify milestone completion, approve or reject milestone releases, and provide reasons for decisions. Auditors ensure quality and compliance before fund release.
    
    \item \textbf{Public/Citizens:} Monitor fund utilization, view project progress, access audit trails, and check impact scores through a read-only public explorer. No authentication required, ensuring complete transparency.
\end{itemize}

\section{Project Objectives}

The primary objectives of this project are:
\begin{enumerate}
    \item To design and implement a blockchain-based system for transparent public fund tracking
    \item To automate fund release through milestone-based escrow smart contracts
    \item To provide immutable audit trails for all transactions
    \item To enable real-time public access to fund utilization data
    \item To implement novel features including NFT approval stamps and public trust index
    \item To demonstrate blockchain's potential in governance applications
\end{enumerate}

\chapter{Motivation}

Traditional public fund management systems face critical challenges that undermine public trust and governance effectiveness. These challenges manifest in multiple dimensions: lack of transparency (opaque allocation, limited public access, delayed reporting), accountability issues (corruption risks, weak audit trails), and manual error-prone processes. These challenges lead to trust deficit, high corruption perception, and economic inefficiency.

\section{Current System Challenges}

Traditional systems operate with limited transparency, making it difficult for citizens to track fund allocation and utilization. Key issues include opaque allocation decisions, restricted public access to fund utilization data, delayed financial reports, and fragmented information across multiple systems.

Accountability mechanisms in traditional systems are weak, leading to corruption risks due to centralized control, weak audit trails that can be altered or deleted, delayed verification processes, and limited real-time oversight capabilities.

\section{Blockchain as a Solution}

Blockchain technology addresses these issues through its inherent properties \cite{nakamoto2008}:

\begin{itemize}
    \item \textbf{Transparency:} All transactions are visible on a public ledger, enabling real-time visibility into fund allocation and utilization, public access to complete transaction history, and immediate availability of financial data.
    
    \item \textbf{Immutability:} Blockchain's cryptographic security ensures tamper-proof audit trails that cannot be altered or deleted, permanent record of all transactions, and cryptographic proof of data integrity.
    
    \item \textbf{Decentralization:} Distributed architecture provides no single point of failure, resistance to censorship, distributed consensus for transaction validation, and reduced dependency on centralized authorities.
    
    \item \textbf{Automation:} Smart contracts enable automated enforcement of rules and regulations, self-executing agreements without intermediaries, reduced need for manual intervention, and consistent application of policies.
\end{itemize}

\section{Expected Benefits}

Implementation of blockchain-based fund tracking is expected to deliver: increased transparency (complete visibility into fund flow), enhanced accountability (immutable records ensure all actions are traceable), reduced corruption (automated processes and transparency reduce opportunities for misappropriation), improved efficiency (automation reduces processing time and administrative costs), and restored public trust (transparent operations rebuild citizen confidence in governance).

International examples include Estonia \cite{estonia2017} (e-governance), Georgia \cite{georgia2016} (land registry), and India's blockchain initiatives \cite{india2021}, demonstrating blockchain's potential in governance.

\chapter{Literature Survey and Existing Systems}

\section{Blockchain Technology}

\begin{definition}[Blockchain]
A distributed ledger technology that maintains a continuously growing list of records (blocks) linked and secured using cryptography \cite{nakamoto2008}. Each block contains a cryptographic hash of the previous block, a timestamp, and transaction data, creating an immutable chain.
\end{definition}

Blockchain technology, introduced through Bitcoin \cite{nakamoto2008}, represents a paradigm shift in how data is stored and verified. Unlike traditional centralized databases, blockchain distributes data across a network of nodes, each maintaining a copy of the entire ledger. This architecture provides several key advantages: decentralization (no single entity controls the network), immutability (once recorded, transactions cannot be altered), transparency (all transactions visible to network participants), and cryptographic security (advanced cryptographic techniques ensure data authenticity).

Smart contracts, first conceptualized by Nick Szabo \cite{szabo1994} in 1994, are self-executing contracts with terms directly written into code. Ethereum \cite{buterin2014}, launched in 2015, brought smart contracts to practical implementation, enabling automated execution of agreements without intermediaries. Smart contracts offer automation (rules enforced automatically), trust (code execution is deterministic and verifiable), efficiency (automated processes reduce processing time and costs), and transparency (contract code is publicly visible).

\section{Blockchain in Governance}

International examples include Estonia \cite{estonia2017} (e-governance, digital identity), Georgia \cite{georgia2016} (land registry), Switzerland \cite{switzerland2019} (e-voting), and UAE \cite{uae2020} (government services). India has initiated blockchain projects for land records, supply chain, and certificate verification \cite{india2021}.

\section{Existing Systems}

\subsection{Traditional Centralized Systems}

Traditional systems for public fund management in India include:

\begin{itemize}
    \item \textbf{PRAGATI (Pro-Active Governance and Timely Implementation) \cite{pragati2015}:} A multi-purpose platform for monitoring and reviewing government programs. While it provides project tracking, it lacks blockchain's immutability and transparency features. The system is centralized, making it vulnerable to data manipulation.
    
    \item \textbf{GeM (Government e-Marketplace) \cite{gem2016}:} An online platform for public procurement. While it digitizes procurement processes, it does not provide immutable audit trails or public transparency. Transactions can be modified, and there is limited public access to utilization data.
    
    \item \textbf{IFMS (Integrated Financial Management System) \cite{ifms2018}:} A financial management system for government departments. It provides accounting and budgeting features but lacks blockchain's transparency and immutability. Audit trails can be altered, and public access is limited.
\end{itemize}

\subsubsection{Limitations of Traditional Systems}

These systems share common limitations:
\begin{itemize}
    \item \textbf{Centralized Architecture:} Single point of control and potential failure
    \item \textbf{Alterable Records:} Data can be modified or deleted, compromising integrity
    \item \textbf{Limited Transparency:} Restricted public access to fund utilization data
    \item \textbf{Manual Processes:} Human intervention required for verification and approval
    \item \textbf{Delayed Reporting:} Financial reports published with significant delays
    \item \textbf{Weak Audit Trails:} Audit records can be tampered with
\end{itemize}

\subsection{Blockchain-Based Systems}

Several blockchain platforms exist, but most focus on different use cases:

\begin{itemize}
    \item \textbf{Stellar \cite{stellar2014}:} A payment network designed for fast, low-cost cross-border transactions. While it uses blockchain technology, it focuses on financial payments rather than public fund tracking and transparency.
    
    \item \textbf{Ripple \cite{ripple2012}:} A payment protocol for financial institutions. Similar to Stellar, it emphasizes payment efficiency rather than governance transparency.
    
    \item \textbf{Hyperledger Fabric \cite{hyperledger2016}:} An enterprise blockchain framework for building permissioned blockchain applications. While suitable for enterprise use, existing implementations focus on supply chain, trade finance, and other business applications rather than public fund transparency.
\end{itemize}

\section{Research Gaps}

Existing systems lack: public transparency, milestone-based automation, immutable audit trails, novel verification mechanisms (NFTs, ZK proofs), and complete stakeholder integration.

\section{Novel Contributions}

This project introduces: complete public transparency, milestone-based escrow, NFT approval stamps, zero-knowledge proofs, public trust index, and comprehensive audit trail.

\chapter{Problem Statement}

\section{Core Problem}

The primary problem is the \textbf{absence of a transparent, accountable, and immutable system for tracking public fund allocation and utilization}. Current systems suffer from: opacity in fund flow, delayed verification and reporting, corruption risks, lack of real-time monitoring, and inadequate audit trails.

\section{Challenges}

Technical challenges include scalability, gas costs, user onboarding, and network dependency. Operational challenges include integration with existing systems, stakeholder adoption, training, and regulatory compliance.

\section{Problem Scope}

This project addresses: transparent fund tracking, automated milestone-based fund release, immutable audit trails, real-time public access, and complete stakeholder integration. Future work includes ERP integration, mobile apps, and cross-chain interoperability.

\chapter{Proposed Design}

\section{System Architecture}

PFTUS follows a layered architecture with four main components: Blockchain Layer (smart contracts), Frontend Layer (web application), Storage Layer (IPFS), and Integration Layer (Web3 libraries). The system connects multiple stakeholder dashboards to blockchain smart contracts, ensuring transparent and accountable fund management.

\begin{figure}[H]
\centering
\includegraphics[width=0.95\textwidth]{figures/system-architecture.png}
\caption{System Architecture Diagram showing the layered structure of PFTUS. The diagram illustrates the Presentation Layer (user dashboards), Application Layer (Next.js frontend), Blockchain Layer (four smart contracts), Storage Layer (IPFS), and Public Explorer (read-only access).}
\label{fig:architecture}
\end{figure}

The blockchain layer consists of four interconnected smart contracts: \textbf{PublicFundProject} (main contract managing projects and milestones), \textbf{ApprovalNFT} (ERC-721 \cite{erc721} compliant contract for minting approval stamps), \textbf{ImpactScoreCalculator} (calculates public trust index), and \textbf{ZKVerifier} (handles privacy-preserving verification).

\section{Smart Contract Design}

\subsection{PublicFundProject Contract}

The main contract manages the complete project lifecycle with the following core functions:

\begin{itemize}
    \item \textbf{Government Functions:}
    \begin{itemize}
        \item \texttt{createProject()}: Initialize new project with milestones, contractor, and auditor
        \item \texttt{allocateFunds()}: Deposit funds into contract escrow
        \item \texttt{pauseProject() / resumeProject()}: Emergency controls
    \end{itemize}
    
    \item \textbf{Contractor Functions:}
    \begin{itemize}
        \item \texttt{requestMilestoneRelease()}: Request funds for next milestone
        \item \texttt{submitProof()}: Submit IPFS hash of completion proof
        \item \texttt{resubmitMilestone()}: Resubmit after rejection
    \end{itemize}
    
    \item \textbf{Auditor Functions:}
    \begin{itemize}
        \item \texttt{approveMilestone()}: Approve milestone and trigger fund release
        \item \texttt{rejectMilestone()}: Reject with reason for correction
    \end{itemize}
    
    \item \textbf{Public View Functions:}
    \begin{itemize}
        \item \texttt{getProjectDetails()}: Retrieve project information
        \item \texttt{getProjectMilestones()}: Get all milestone data
        \item \texttt{getAuditTrail()}: Complete transaction history
        \item \texttt{getProjectProgress()}: Calculate completion percentage
    \end{itemize}
\end{itemize}

\subsection{Supporting Contracts}

\textbf{ApprovalNFT Contract:} ERC-721 \cite{erc721} compliant contract for minting soulbound (non-transferable) approval stamps. Each milestone approval automatically mints an NFT providing immutable proof of verification.

\textbf{ImpactScoreCalculator Contract:} Calculates public trust index using weighted formula:
\begin{equation}
\text{Score} = 0.4 \times T_{\text{on-time}} + 0.3 \times B_{\text{adherence}} + 0.2 \times A_{\text{rate}} + 0.1 \times T_{\text{transparency}}
\end{equation}
where each component is scored 0-100, resulting in a final score of 0-100.

\textbf{ZKVerifier Contract:} Simplified zero-knowledge proof verification for privacy-preserving document verification, allowing citizens to verify spending without accessing sensitive invoices.

\section{Workflow Design}

\subsection{Project Lifecycle Workflow}

The complete project lifecycle follows a sequential workflow ensuring transparency and accountability at each stage. The workflow consists of twelve distinct phases:

\begin{enumerate}
    \item \textbf{Project Creation:} Government creates project with defined milestones, assigns contractor and auditor addresses, and sets budget and timeline. Project status is set to PENDING.
    
    \item \textbf{Fund Allocation:} Government deposits funds into the smart contract escrow. Funds are locked and cannot be withdrawn until milestones are approved. Project status changes to ACTIVE.
    
    \item \textbf{Milestone Request:} Contractor requests release for a specific milestone. The milestone status changes from PENDING to IN\_PROGRESS. Sequential enforcement ensures previous milestones are completed first.
    
    \item \textbf{Proof Submission:} Contractor uploads completion proof documents to IPFS and submits the IPFS hash to the contract. Milestone status changes to UNDER\_REVIEW.
    
    \item \textbf{Audit Review:} Auditor reviews the submitted proof and makes a decision: Approve (triggers automatic fund release, NFT minting, and score update) or Reject (allows contractor to resubmit with corrections).
    
    \item \textbf{Fund Release:} Upon approval, funds are automatically transferred to the contractor's address. This is an atomic operation that cannot be reversed.
    
    \item \textbf{NFT Minting:} An approval NFT is automatically minted, providing immutable proof of verification. The NFT is soulbound (non-transferable) and linked to the milestone.
    
    \item \textbf{Impact Score Update:} The public trust index is recalculated based on completion time, budget adherence, approval rate, and transparency metrics.
    
    \item \textbf{Public Viewing:} Citizens can view project progress, fund utilization, complete audit trail, and NFT approvals in real-time through the Public Explorer.
    
    \item \textbf{Next Milestone:} If additional milestones exist, the process repeats from step 3, maintaining sequential order.
    
    \item \textbf{Project Completion:} When all milestones are approved, the project status changes to COMPLETED, and the final impact score is calculated.
\end{enumerate}

\subsection{Milestone Processing Flow}

Each milestone follows a specific state machine with multiple verification steps. Milestones progress through states: PENDING → IN\_PROGRESS → UNDER\_REVIEW → APPROVED or REJECTED. Upon approval, automatic actions include fund release, NFT minting, and impact score update. Rejected milestones can be resubmitted after corrections. Sequential enforcement ensures milestone N can only be requested after Milestone N-1 is approved.

\begin{figure}[H]
\centering
\includegraphics[width=0.95\textwidth]{figures/project-lifecycle.png}
\caption{Complete Project Lifecycle Workflow showing the 12-step process from project creation to completion}
\label{fig:workflow}
\end{figure}

\begin{figure}[H]
\centering
\includegraphics[width=0.9\textwidth]{figures/milestone-flow.png}
\caption{Milestone Status Flow Diagram showing state transitions from PENDING through IN\_PROGRESS, UNDER\_REVIEW, to APPROVED or REJECTED}
\label{fig:milestone-flow}
\end{figure}

\section{Security Design}

Role-based access control (Government, Contractor, Auditor), reentrancy protection using OpenZeppelin's \texttt{ReentrancyGuard} \cite{openzeppelin}, pausable mechanism, input validation, sequential milestone release, and immutable audit trail.

\section{Novel Features}

Milestone-based escrow (funds locked until verified completion), NFT approval stamps (ERC-721 \cite{erc721} soulbound tokens), zero-knowledge proofs (privacy-preserving verification), public trust index (automated scoring system), and complete audit trail (immutable record of all transactions).

\chapter{Implementation}

\section{Technology Stack}

\textbf{Blockchain:} Solidity 0.8.20, Hardhat \cite{hardhat}, OpenZeppelin v5.0 \cite{openzeppelin}, Polygon Mumbai testnet \cite{polygon}. \textbf{Frontend:} Next.js 14 \cite{nextjs}, React 18, TypeScript, Tailwind CSS, ethers.js v6 \cite{ethers}. \textbf{Storage:} IPFS \cite{ipfs} (Pinata), MetaMask \cite{metamask} for wallet integration.

\section{Smart Contract Implementation}

\subsection{PublicFundProject Contract Structure}

The main \texttt{PublicFundProject} contract (600+ lines of Solidity) implements the core functionality of the system. The contract is built using OpenZeppelin's security libraries \cite{openzeppelin}, including \texttt{ReentrancyGuard} for protection against reentrancy attacks, \texttt{Pausable} for emergency controls, and \texttt{AccessControl} for role-based permissions.

The contract implements milestone-based escrow (funds locked until verified completion), sequential milestone processing (enforces order of milestone completion), comprehensive audit trail (every action recorded with timestamp, actor, and details), automatic fund release (funds released automatically upon approval), and integration with supporting contracts (seamless interaction with ApprovalNFT and ImpactScoreCalculator).

\subsection{Fund Release Conditions}

Funds are released only when all of the following conditions are met:
\begin{enumerate}
    \item Milestone has been requested by the contractor
    \item Proof documents have been submitted (IPFS hash stored on-chain)
    \item Milestone has been approved by the assigned auditor
    \item Previous milestones (if any) have been completed and approved
    \item Contract is not paused
    \item Sufficient funds are available in the contract escrow
\end{enumerate}

This multi-condition approach ensures that funds are only released for verified, completed work, preventing premature or unauthorized fund access.

\section{Frontend Implementation}

\subsection{Dashboard Architecture}

The frontend provides four specialized dashboards, each tailored to specific stakeholder needs:

\begin{itemize}
    \item \textbf{Government Dashboard:} Create projects, allocate funds, monitor all projects with status indicators and progress bars, pause/resume projects in emergency situations, and view aggregated statistics.
    
    \item \textbf{Contractor Dashboard:} View assigned projects (automatically filtered by connected wallet address), request milestones with status indicators, submit proofs by uploading documents to IPFS, track milestone approval/rejection status in real-time, and view payment history.
    
    \item \textbf{Auditor Dashboard:} Review pending milestones with proof document links, approve/reject milestones with reasons, access IPFS-hosted proof documents for verification, track all approvals and rejections with timestamps, and view project details and milestone progress.
    
    \item \textbf{Public Explorer:} Read-only access to all projects with search and filter capabilities, visual representation of fund flow, complete transaction history for each project, impact scores with detailed breakdowns, NFT approval stamps for verified milestones, and comprehensive project information. No authentication required, ensuring complete transparency.
\end{itemize}

\subsection{Web3 Integration}

The frontend uses ethers.js v6 \cite{ethers} for blockchain interaction, providing provider management (automatic detection and connection to MetaMask \cite{metamask} wallet), contract interaction (type-safe contract function calls), transaction handling (asynchronous transaction submission with confirmation waiting), event listening (real-time event listening for blockchain state changes), network management (automatic network switching to correct blockchain), and comprehensive error handling.

Proof documents are uploaded to IPFS \cite{ipfs} using the Pinata service for decentralized storage. Files are uploaded to IPFS through Pinata API, returning IPFS hash (CID), which is then stored on-chain in smart contracts, ensuring document integrity while keeping blockchain storage costs low.

\section{Deployment}

Contracts deployed to local Hardhat network and Polygon Mumbai testnet using Hardhat deployment scripts.

\chapter{Analysis}

\section{Testing and Results}

\subsection{Smart Contract Testing}

Smart contracts were thoroughly tested across multiple dimensions:

\begin{itemize}
    \item \textbf{Functionality Testing:} Verified that projects are created with correct parameters, funds are correctly deposited into escrow, milestone processing flows work correctly, funds are automatically transferred upon approval, and sequential milestone ordering is enforced.
    
    \item \textbf{Security Testing:} Tested with malicious contracts attempting reentrancy during fund transfers (all attempts blocked by \texttt{ReentrancyGuard}), verified that unauthorized users cannot perform restricted actions, tested with invalid inputs to ensure proper validation, and tested boundary conditions.
    
    \item \textbf{Gas Optimization Testing:} Measured gas costs for all major operations. Create Project: ~500,000 gas (depends on number of milestones), Allocate Funds: ~100,000 gas, Request Milestone: ~80,000 gas, Submit Proof: ~60,000 gas, Approve Milestone: ~200,000 gas (includes fund transfer and NFT minting), Reject Milestone: ~50,000 gas, Mint NFT: ~150,000 gas.
\end{itemize}

On Polygon Mumbai testnet, these costs translate to minimal fees (less than \$0.01 per transaction), making the system economically viable. However, on Ethereum mainnet, gas costs would be significantly higher, requiring optimization or Layer 2 solutions.

\subsection{Frontend Testing}

Frontend testing covered UI usability (verified intuitive dashboards), responsiveness (tested on multiple screen sizes), Web3 integration (wallet connection, network switching, transaction handling), error handling (network failures, transaction rejections, invalid inputs), and state synchronization (UI updates correctly after blockchain state changes).

\section{Key Achievements}

\textbf{Transparency:} All transactions visible on blockchain with complete audit trail. \textbf{Accountability:} Automated fund release, role-based access control, NFT proof stamps. \textbf{Security:} Reentrancy protection, access control, input validation, emergency pause.

\section{Limitations}

Current limitations: High gas costs on mainnet, network dependency, user onboarding requirements, simplified ZK implementation. Challenges faced: Contract complexity, frontend-blockchain synchronization, comprehensive error handling.

\section{Comparison with Traditional Systems}

A comprehensive comparison between PFTUS and traditional systems:

\begin{itemize}
    \item \textbf{Transparency:} PFTUS provides complete transparency with all transactions visible on public ledger, while traditional systems offer limited transparency with restricted access and delayed reports.
    
    \item \textbf{Immutability:} PFTUS ensures tamper-proof records that cannot be altered, while traditional systems have alterable records that can be modified or deleted.
    
    \item \textbf{Updates:} PFTUS provides real-time updates with immediate blockchain updates, while traditional systems have delayed updates with periodic reports often months late.
    
    \item \textbf{Public Access:} PFTUS offers open public access with no authentication required for viewing, while traditional systems have restricted access with limited public access.
    
    \item \textbf{Audit Trail:} PFTUS provides automated immutable blockchain records, while traditional systems have manual alterable records with weak trails.
    
    \item \textbf{Fund Release:} PFTUS automates fund release through smart contract enforcement, while traditional systems require manual human approval.
    
    \item \textbf{Cost:} PFTUS has low cost per transaction (testnet costs minimal, Layer 2 solutions available), while traditional systems have high costs with administrative overhead and manual processing.
    
    \item \textbf{Security:} PFTUS provides high security with cryptographic security and decentralization, while traditional systems have moderate security with centralized architecture vulnerable to attacks.
    
    \item \textbf{Corruption Risk:} PFTUS has low corruption risk with automated processes and transparency, while traditional systems have high corruption risk with manual processes and limited oversight.
\end{itemize}

\chapter{Conclusion}

\section{Summary}

This project successfully demonstrates blockchain technology's application to public fund management, addressing transparency, accountability, and corruption issues. PFTUS provides: transparent fund tracking, automated accountability, immutable audit trail, public trust through independent verification, and novel features (NFT stamps, impact scores, ZK proofs).

\section{Contributions}

This project makes several significant contributions:

\subsection{System Design Contributions}

\begin{itemize}
    \item \textbf{Complete Blockchain-Based System:} Designed and implemented a comprehensive blockchain-based system for public fund tracking that addresses real-world governance challenges.
    
    \item \textbf{Multi-Stakeholder Architecture:} Created an architecture that seamlessly integrates four distinct stakeholder roles (Government, Contractor, Auditor, Public) with appropriate access controls and interfaces.
    
    \item \textbf{Automated Workflow:} Implemented automated workflow that reduces manual intervention while maintaining security and accountability.
\end{itemize}

\subsection{Technical Contributions}

\begin{itemize}
    \item \textbf{NFT Approval Stamps:} Novel use of ERC-721 \cite{erc721} tokens as soulbound (non-transferable) approval stamps, providing immutable proof of verification.
    
    \item \textbf{Public Trust Index:} Automated scoring system that calculates project performance based on multiple weighted metrics, providing real-time transparency indicators.
    
    \item \textbf{Simplified ZK Proofs:} Implemented zero-knowledge proof verification for privacy-preserving document verification, balancing transparency with privacy.
    
    \item \textbf{Gas Optimization:} Applied multiple gas optimization techniques to ensure cost-effective operations on blockchain networks.
\end{itemize}

\subsection{Practical Contributions}

\begin{itemize}
    \item \textbf{Proof of Concept:} Demonstrated that blockchain technology can be effectively applied to public fund management, providing a working prototype for real-world deployment.
    
    \item \textbf{User Experience:} Created user-friendly interfaces that make blockchain interactions accessible to non-technical users.
    
    \item \textbf{Integration Patterns:} Established patterns for integrating blockchain smart contracts with traditional web applications and decentralized storage.
\end{itemize}

\section{Future Work}

Short-term: Mobile applications, advanced analytics, notification system, multi-currency support, enhanced ZK proofs. Long-term: Government ERP integration, AI-based anomaly detection, cross-chain interoperability, decentralized governance (DAO), and Layer 2 scalability solutions \cite{polygon}.

\section{Final Remarks}

Blockchain technology holds immense potential for transforming governance. This project demonstrates that blockchain can provide the transparency, accountability, and trust that modern governance systems require. PFTUS serves as a proof-of-concept that can be extended for real-world deployment. As noted in various studies on blockchain in governance \cite{blockchain-gov}, the technology's potential is only beginning to be realized.

\clearpage

% Bibliography
\begin{thebibliography}{99}

\bibitem{nakamoto2008}
Nakamoto, S. (2008). Bitcoin: A Peer-to-Peer Electronic Cash System. Available at: \url{https://bitcoin.org/bitcoin.pdf}

\bibitem{buterin2014}
Buterin, V. (2014). A Next-Generation Smart Contract and Decentralized Application Platform. Ethereum White Paper. Available at: \url{https://ethereum.org/en/whitepaper/}

\bibitem{szabo1994}
Szabo, N. (1994). Smart Contracts. Available at: \url{http://www.fon.hum.uva.nl/rob/Courses/InformationInSpeech/CDROM/Literature/LOTwinterschool2006/szabo.best.vwh.net/smart.contracts.html}

\bibitem{estonia2017}
Estonia e-Governance Academy. (2017). Estonian e-Government Ecosystem: Foundation, Applications, Outcomes. Available at: \url{https://ega.ee/wp-content/uploads/2017/05/Estonian-eGovernment-ecosystem.pdf}

\bibitem{georgia2016}
Government of Georgia. (2016). Blockchain-based Land Registry System. Available at: \url{https://www.napier.gov.ge/en/news/view/156}

\bibitem{switzerland2019}
Swiss Federal Chancellery. (2019). E-Voting System Using Blockchain Technology. Available at: \url{https://www.bk.admin.ch/bk/en/home/political-rights/e-voting.html}

\bibitem{uae2020}
UAE Government. (2020). UAE Blockchain Strategy 2021. Available at: \url{https://u.ae/en/about-the-uae/strategies-initiatives-and-awards/federal-governments-strategies-and-plans/uae-blockchain-strategy-2021}

\bibitem{india2021}
NITI Aayog. (2021). Blockchain: The India Strategy. Available at: \url{https://www.niti.gov.in/blockchain-india-strategy}

\bibitem{pragati2015}
Government of India. (2015). PRAGATI: Pro-Active Governance and Timely Implementation. Available at: \url{https://pragati.gov.in/}

\bibitem{gem2016}
Government e-Marketplace. (2016). GeM: Government e-Marketplace Portal. Available at: \url{https://gem.gov.in/}

\bibitem{ifms2018}
Government of India. (2018). Integrated Financial Management System (IFMS). Available at: \url{https://ifms.nic.in/}

\bibitem{stellar2014}
Stellar Development Foundation. (2014). Stellar Consensus Protocol. Available at: \url{https://www.stellar.org/developers/guides/concepts/scp.html}

\bibitem{ripple2012}
Ripple Labs. (2012). Ripple Protocol Consensus Algorithm. Available at: \url{https://ripple.com/files/ripple_consensus_whitepaper.pdf}

\bibitem{hyperledger2016}
Hyperledger Foundation. (2016). Hyperledger Fabric: A Distributed Operating System for Permissioned Blockchains. Available at: \url{https://www.hyperledger.org/use/fabric}

\bibitem{openzeppelin}
OpenZeppelin. (2023). OpenZeppelin Contracts: Secure Smart Contract Library. \url{https://docs.openzeppelin.com/contracts}

\bibitem{polygon}
Polygon Technology. (2023). Polygon Documentation. \url{https://docs.polygon.technology}

\bibitem{hardhat}
Nomic Foundation. (2023). Hardhat Documentation. \url{https://hardhat.org/docs}

\bibitem{nextjs}
Vercel. (2023). Next.js Documentation. \url{https://nextjs.org/docs}

\bibitem{ethers}
ethers.js. (2023). ethers.js Documentation. \url{https://docs.ethers.org}

\bibitem{ipfs}
Protocol Labs. (2023). IPFS Documentation. \url{https://docs.ipfs.tech}

\bibitem{metamask}
MetaMask. (2023). MetaMask Documentation. \url{https://docs.metamask.io/}

\bibitem{erc721}
Ethereum Improvement Proposal 721. (2018). ERC-721 Non-Fungible Token Standard. Available at: \url{https://eips.ethereum.org/EIPS/eip-721}

\bibitem{blockchain-gov}
Kumar, R., et al. (2020). Blockchain Applications in Government: A Comprehensive Survey. International Journal of Information Management, 52, 102-115.

\end{thebibliography}

\end{document}
